% !TeX root = ../教材题库_全部条目_转写工作稿.tex
% Source: XJ-B1-C06-PR-P230-Y1031
% 第6章 / 6.3 统计图表 / 教材 p.230 / PDF 第237页
% OCR未校对：公式与文字均可能有误；未完成数学TeX转写；禁止直接发给学生
% 原题及相邻原解截图（核对用）：../题目原图/XJ-B1-C06-PR-P230-Y1031-01.webp
% 本文件只在转写完成后作为题干使用；OCR 条目尚不是可靠题干。
\begin{bankitem}{p.230 练习 第3题}{第6章 / 6.3 统计图表 / 教材 p.230 / PDF 第237页}{XJ-B1-C06-PR-P230-Y1031}{OCR未校对：公式与文字均可能有误；未完成数学TeX转写；禁止直接发给学生}
\begin{ocrdraft}
[教材p.230 / PDF第237页]
\par
3．
\par
一种袋装食品用生产线自动装填，每袋质量大约为50g，但由于某些原
\par
因，每袋食品不会恰好是50g.现随机抽取100袋食品，测得的质量数据如下：
\par
单位：g
\par
57
\par
46
\par
49
\par
54
\par
55
\par
58
\par
49
\par
51
\par
49
\par
52
\par
54
\par
51
\par
55
\par
60
\par
56
\par
47
\par
47
\par
48
\par
53
\par
50
\par
52
\par
40
\par
45
\par
53
\par
52
\par
46
\par
48
\par
47
\par
53
\par
47
\par
53
\par
44
\par
47
\par
50
\par
52
\par
53
\par
47
\par
45
\par
48
\par
54
\par
52
\par
48
\par
46
\par
49
\par
52
\par
59
\par
53
\par
50
\par
43
\par
53
\par
46
\par
49
\par
49
\par
44
\par
52
\par
42
\par
49
\par
43
\par
47
\par
46
\par
48
\par
59
\par
45
\par
45
\par
46
\par
52
\par
55
\par
47
\par
49
\par
50
\par
54
\par
47
\par
48
\par
44
\par
47
\par
53
\par
58
\par
52
\par
48
\par
53
\par
49
\par
56
\par
56
\par
53
\par
41
\par
48
\par
（D为了获得样本数据的分布情况，试制作频率分布表；
\par
（2）绘制频率分布直方图及频率分布折线图．
\end{ocrdraft}
\end{bankitem}
